\section{Procedura}

\begin{enumerate}
  \item Durante il meeting, un membro del team designato prende appunti;
  \item Al termine del meeting, gli appunti vengono strutturati nel formato descritto nella Sezione precedente;
  \item Il verbale viene scritto in LaTeX nella cartella \texttt{PB/verbali/};
  \item Viene aperto un branch \texttt{docs/verbale-YYYY-MM-DD} contenente solo il nuovo verbale;
  \item Viene aperta una pull request che viene mergiata dopo verifica da parte di un altro partecipante al meeting (per garantire l'accuratezza del contenuto);
  \item Il verbale è quindi disponibile a tutto il team.
\end{enumerate}

\subsection{Frequenza}

Si raccomanda di redigere un verbale per:

\begin{itemize}
  \item Ogni meeting di pianificazione settimanale;
  \item Ogni decisione strategica significativa (es. approvazione di una nuova strategia per l'ensemble core);
  \item Ogni meeting con stakeholder esterni (es. presentazione preliminare a Fund of Funds).
\end{itemize}

Riunioni puramente operative di breve durata (es. stand-up giornalieri) possono essere documentate in modo più sintetico o via messaggistica del team.
